\section{Estrategia de partición del espacio de claves}\label{sec:partition}
El espacio de búsqueda es el conjunto de todas las cadenas de longitud fija sobre el alfabeto base-36. Para seis caracteres su tamaño es $N = 36^6 = 2{,}176{,}782{,}336$. Cada índice entero $i \in [0, N)$ se convierte determinísticamente en una clave mediante codificación posicional, por lo que un intervalo de índices corresponde a un conjunto de claves distintas y dos intervalos disjuntos no se traslapan. El main divide $[0,N)$ en bloques de tamaño fijo $C$, obteniendo $\lceil N/C \rceil$ bloques (por ejemplo, $C = 5\times10^6$ da 436 bloques) servidos desde una queue.

La asignación es dinámica. El main conserva una queue de bloques pendientes y un registro de los asignados (\emph{inflight}). Ante cada \texttt{REQUEST}, si hay pendientes, entrega el siguiente mediante \texttt{ASSIGN}$(c,\, cC,\, \min((c{+}1)C, N))$ y lo marca como asignado. Si no quedan pendientes pero sí asignados, responde \texttt{NOWORK} porque alguno podría reencolarse, y si no queda ninguno, marca el espacio como agotado y difunde \texttt{STOP}. Un bloque solo se da por terminado cuando llega su \texttt{DONE}.
